\SourceLeaf{413}%
\par
Als Abszissenachse wähle ich die Gerade \(\mathrm{(II)}\), die \((4),(5)\) in ihrem Mittelpunkt \((6)\) im rechten Winkel schneidet. Ich zähle die Abszissen von \((6)\) aus auf der Seite von \((4),(5)\) positiv, auf der der Winkel von \((4)\) nach \((5)\) kleiner als \(180^0\) erscheint, also auf der rechten Seite, wenn man die Winkel von links nach rechts wachsen lässt. Die Ordinaten seien in der Richtung von \((6)\) nach \((5)\) positiv. Auf \(\mathrm{(II)}\) bezeichne ich die Punkte mit den Abszissen
\[c.\operatorname{cotang}\vartheta=n-a,\quad c.\operatorname{cotang}\vartheta'=n,\quad c.\operatorname{cotang}\vartheta''=n+b\]
mit \((1^*),(2^*),(3^*)\). Das sind die Mittelpunkte der drei Kreise, die jeweils durch \((1),(2),(3)\) und zugleich alle durch \((4)\) und \((5)\) gehen. Die Radien dieser Kreise sind
\[\frac c{\sin\vartheta}=\BareRoot{(cc+(n-a)^2)},\quad\frac c{\sin\vartheta'}=\BareRoot{(cc+nn)},\quad\frac c{\sin\vartheta''}=\BareRoot{(cc+(n+b)^2)}\]
oder, wenn man \(\frac c{\sin\vartheta'}=r\) setzt, ergeben sich für die beiden anderen \(\BareRoot{(rr-2an+aa)},\BareRoot{(rr+2bn+bb)}\). Schließlich sei \((7)\) der Schnittpunkt von \(\mathrm{(I)}\) und \(\mathrm{(II)}\). \(T\) und \(t\) seien die Abstände der Punkte \((2)\) und \((2^*)\) von \((7)\), und \(\varphi\) sei der Winkel zwischen den gleichnamigen Strahlen dieser Geraden, gemessen von \(\mathrm{(I)}\) nach \(\mathrm{(II)}\) und im gewählten Drehsinn positiv. Somit ist die Abszisse von \((7)=n-t\), und damit lauten die Koordinaten der drei Beobachtungsorte
\[\begin{array}{rcll}(1)&.\ .\ .\ .\ .&n-t+(T-A).\cos\varphi,&(T-A).\sin\varphi\\
(2)&.\ .\ .\ .\ .&n-t+T.\cos\varphi,&T.\sin\varphi\\
(3)&.\ .\ .\ .\ .&n-t+(T+B)\cos\varphi,&(T+B).\sin\varphi\end{array}\]
\par
Die drei unbekannten Größen \(t,T,\varphi\) lassen sich aus den folgenden Gleichungen bestimmen, wobei \(x\) zur Abkürzung für \(\cos\varphi\) steht.
\[\begin{aligned}tt+TT-2tTx&=rr\quad.\ .\ .\ .\ .\ .\ [1]\\
(t-a)^2+(T-A)^2-2(t-a)(T-A)x&=rr-2an+aa\\
(t+b)^2+(T+B)^2-2(t+b)(T+B)x&=rr+2bn+bb\end{aligned}\]
\par
Anstelle der beiden letzten Gleichungen verwende ich die folgenden, die sich ergeben, wenn man die beiden letzten Gleichungen jeweils von der ersten abzieht.
\[\begin{aligned}2at+2AT-2an-AA&=(2At+2aT-2aA)x\quad.\ .\ .\ [2]\\
2bt+2BT-2bn+BB&=(2Bt+2bT+2bB)x\quad.\ .\ .\ [3]\end{aligned}\]
